% Transcription — fonds Alexandre Grothendieck, shelfmark 37, batch 2 (pages 21-40).
% Established from the facsimile archives/batches/37/batch-02.pdf (a copy of
% 37.pdf fetched through the project's relay, grothendieck-relay.onrender.com;
% PDF page k = archivists' page 20+k, no cover sheet), per the skill
% transcribe-grothendieck. The folder is 99 pages in five batches (1-20, 21-40,
% 41-60, 61-80, 81-99); this is the second, done in parallel with the others.
%
% Pass: Opus 5.5 (claude-opus-5-5), 2026-09-26 — first pass, unchecked against the pages by a human.
%
% Pages skipped: 28 and 31 (blank cream sheets), 36 (a blank IHES letterhead
% leaf, upside down, nothing written on it).
% Pages noted only (not re-set): 26, 33, 34.
% Pages transcribed: 21-25, 27, 29, 30, 32, 35, 37-40.
%
% Decisions, item by item, with the evidence:
% - Pages 21, 23, 24, 25, 27, 32, 35, 37: his own manuscript notes (IHES
%   letterhead on 21, 23, 32, 37; a half-sheet slip on 27), lists of things to
%   do for the re-edition of SGA 4, exposé by exposé. Transcribed in full.
%   Page 21 is not the continuation of his letter to Verdier of 18.5.1971
%   (batch 1, p. 14): it is a list headed « SGA 4 VI », with no salutation,
%   closing or letter text. Page 24 is a manuscript « Note pour la page 8 »
%   of SGA 4 IX (heading « SGA 4 IX p. 42 »), the note to which the published
%   VI 1.27 refers (« cf. IX, note p. 42 »); being in his hand, it is given
%   entire, including the block he cancelled by diagonal strokes.
% - Page 22: a typed leaf headed in pencil « SGA 4 VI 2 », an « Exercice 2. »
%   (pencil numeral over a blank): for a small category C with a Karoubi
%   envelope C', conditions (i)-(iv) for Ind(C) to be a U-topos, with E_PF.
%   DECISION: his own UNPUBLISHED typescript, transcribed in full with his ink
%   and pencil changes. Evidence that it is not the published text: the
%   Orgogozo re-edition of SGA 4 t. 2 (normalesup.org/~forgogozo/SGA4, tome2,
%   consulted 2026-09-27) has in VI only Exercices 1.28, 2.16, 2.17, 2.18,
%   3.11, 3.12 ..., none on Ind(C) as a topos, and no « enveloppe de Karoubi »
%   in Exposé VI; the related published statements are VI 1.24-1.25 and the
%   definition of « parfait » 2.9.1. Exposé VI is his (« par A. Grothendieck »
%   in the published volume), and the leaf's pencil label is in the same hand
%   as his notes of pp. 21, 23.
% - Page 26: a typed leaf, upside down, paginated « 19 » (and a pencil number
%   top right, read 92 with doubt), of SGA 7 Exposé VII (biextensions),
%   2.1.1-2.2.3: published text (SGA 7 I, LNM 288). One French \note with his
%   marks: two long pencil diagonals cancelling the leaf, a circled psi in
%   the left margin, a pencil bracket round (2.1.1.2), and in blue ink the
%   heading « Sections unité. » inserted before 2.2, « De plus, » over a
%   struck word, « donc » inserted.
% - Pages 29-30: a typed letter in the first person, addressed to a
%   correspondent he calls « tu », on how to present « 2.1.1 » (a statement on
%   champs and ind-L-champs, which the correspondent's text proves by the
%   proper base change theorem) « en réédition de SGA 4 », with the enclosed
%   « texte joint » (p. 30: a Théorème 2.1.1 and Corollaire 2.1.2 as he
%   proposes them). DECISION: his own letter and his own unpublished draft
%   statement, transcribed in full with every ink change. Evidence: the writer
%   announces « Je pense présenter les choses ainsi en réédition de SGA 4 » —
%   the organiser of that re-edition (cf. batch 1) — and the ink additions
%   and the closing handwritten remark of p. 30 (« En tout, ça devrait faire 3
%   pages au Chap. VII au maximum ... ») are in his hand, the one of pp. 21,
%   23; the typed text says « texte joint » and p. 30 is typed on the same
%   machine and cancels « champ » for « ind-L-champ » as the letter asks. The
%   correspondent's text is divided into chapters (« Chap. VII », « Chap. I »)
%   and is not named; the recipient is not guessed. No date on the leaves.
% - Pages 33-34: a typed letter in English (no heading, no signature, no
%   date) answering his queries on SGA 4 exposés XIII, XIX, X, XII and IX; the
%   writer calls « my paper on alg. approx » M. Artin's « Algebraic
%   approximation of structures over complete local rings », so it is Artin's
%   (and p. 32, in his hand, lists « Artin SGA 4 » with the same points). Another
%   person's letter: one \note per page, not re-set; his marks (a red ink box
%   round the exp XIX paragraph, red and pencil strokes, pencil separating
%   lines) given in full. The hand-inked l (ell) of p. 34 is the typist's.
% - Pages 38-39: a typed text headed in pencil « SGA 4 VIII »: « 10.
%   Compléments : espaces algébriques et étendues algébriques », 10.1-10.5
%   (10.5 empty). DECISION: an UNPUBLISHED typescript, transcribed in full
%   with his marks. Evidence: the published Exposé VIII (re-edition t. 2,
%   consulted 2026-09-27) has nos 1-9 only, and so does the original IHES
%   typescript of VIII (orgogozo.perso.math.cnrs.fr, SGA4-VIIIo.pdf), which
%   ends at 9.4; no « étendue algébrique » occurs in either. The two
%   handwritten entries of the terminology list (batch 3, p. 44) cite
%   « étendue algébrique VIII 10.2 » and « espace algébrique VIII 10.4 »,
%   i.e. this section. Typist not established; the ink and pencil marks and
%   the handwritten footnote on Artin are his, as elsewhere presumable
%   rather than proven.
% - Page 40: the first typed leaf of the bilingual list of SGA 4 terms and
%   notations, under his ink heading « Terminologie, notations SGA 4 Vol. 2 /
%   Attention, faire deux listes disjointes ! ». Transcribed in full with every
%   mark, exactly as batch 3 treats its continuation (pp. 42-44), so that the
%   list reads on across the two files: same three columns separated by
%   dashes, « (entouré) » for the notations he rings in pencil.
%
% His pagination and numbered runs: none on his manuscript notes. The typed
% SGA 7 leaf (p. 26) is numbered « 19 ». The letter (p. 29) opens « VII 2 »
% (the correspondent's chapter and section). The exercise of p. 22 is
% « Exercice 2. », conditions (i), (i bis), (i ter), (ii), (iii), (iv); the
% VIII typescript runs 10, 10.1-10.5. Dates: none in the batch.
%
% What the batch is about. Preparing the re-edition of SGA 4: lists of
% corrections by exposé (VI: introduction, E/phi to phi\E, the exercise on
% finite topoi, « exists lim-> finies »; VII: 5.6 a) X~_et locally coherent,
% coherent iff X is, 5.6 b) X~_et = 2-lim of X~_i for a filtered projective
% system with affine transition maps; IX p. 42; XII, XIII 3.1, XVI 1, XVII
% 4.2.3, XVIII, XIX; a preface to the second edition, Deligne's
% contributions, Verdier's C.D.), Artin's answers and the list drawn from
% them, the note to IX p. 8 proving that constructibility (IX 2.3) agrees
% with VI 1.9.3 and that X~_et satisfies VI 1.25, an unpublished exercise on
% Ind(C) as a topos, a letter proposing to state « 2.1.1 » (equivalence
% G(X) -> G'(X') for ind-L-champs under universally locally L-1-aspheric
% base change) within the re-edition, a planned VIII §10 on algebraic spaces
% and « étendues algébriques », and the start of the terminology list.
%
% Notation fixed once: X~_et as \widetilde{X}_{\text{ét}} (the typescripts
% write both « et » and « ét »); lim-> and lim<- as \varinjlim, \varprojlim;
% the underlined O of the typescripts as \underline{O}; \mathbb{U} for the
% underlined U. Underlined words in running text are \emph.
%
% Hardest pages: 21, 25, 32 and 37 (fast notes, much of the prose lost), 24
% (a cancelled half-page under diagonals). Apparatus: 63 \ill{}, 27
% \uncertain{} (counted on the finished file).

\documentclass[11pt,a4paper]{article}
\input{../preamble/grothendieck}

\folder{37}
\batch{2}
\pages{21}{40}
\dating{1967-1971}
\watermark{Édition de démonstration}
\foldertitle{Rééditions SGA, EGA [SGA 4 à SGA 7, EGA I et EGA II] : tapuscrits et copies de tapuscrit annoté (s.d.), lettres (1967-1971, s.d.), notes manuscrites (s.d.)}

\begin{document}

\section*{Notes pour la réédition de SGA 4}

\page{21}

\note{à l'encre bleue, sur papier à en-tête de l'IHES ; écriture rapide}

SGA 4 VI

Introduction : \uncertain{signaler} \ill{} « \uncertain{sans doute} lié au
fait que \ill{} cond. d'un \ill{} et un \ill{}. »

§4 \quad changer $E_{/\varphi}$ en $\varphi\backslash E$

§3 (?) \quad Changer l'exercice sur les topos finis.

§4 \quad \ill{} \uncertain{vif} au cas \ill{} \ill{} mieux dans l'exercice de
IV 9.1 ?

§1 \quad 1.25 \ldots\ \ill{} \uncertain{covariant}.
Remplacer « $\exists$ sommes \struck{\ill{}} » \add{\ill{}} et « \ill{} » par
« $\exists\ \varinjlim$ finies »\note{la ligne insérée au-dessus du mot biffé
n'est pas lue}

§ , \quad Exercice : $\varinjlim$ topos (\ill{} \uncertain{trivialement})

\page{22}

\note{feuillet dactylographié, coté au crayon en haut à gauche « SGA 4 VI 2 »
(le chiffre romain encadré). Texte non publié (voir l'en-tête) : on le donne
entier. Les mots tapés dans l'interligne sont donnés sans marque ; ses ajouts
à l'encre et au crayon sont en \add{}. Les biffures faites à la machine ne sont
pas relevées}

\emph{Exercice} \add{2}. \quad Soit \add{en} C une \add{petite} catégorie,
\add{($C'$ une enveloppe de Karoubi de $C$ (IV 7.5 b)))}. Montrer que les
conditions suivantes sont équivalentes :\note{« 2 » au crayon dans un blanc
de la machine ; « en » à l'encre au-dessus de « Soit » ; « petite » tapé
au-dessus d'un mot noirci ; la parenthèse sur $C'$ est écrite à l'encre dans
la marge de droite et appelée par un trait après « catégorie, »}

(i) La catégorie $\mathrm{Ind}(C)$ (I 8.2) est un $\mathbb{U}$-topos.

(i bis) Le foncteur canonique $\mathrm{Ind}(C) \to \widehat{C}$ admet un adjoint
à droite qui est exact à gauche.

(i ter) Il existe une topologie sur $C$ pour laquelle les faisceaux soient les
préfaisceaux ind-représentables.

\struck{(ii) Dans $C'$, les limites inductives finies sont représentables, il
existe sur $C'$ une topologie dont les familles $(X_i \to X)_{i \in I}$ sont
celles qui sont raffinées par une famille épimorphique finie}\note{premier état
de (ii), encadré au crayon et barré de quatre traits obliques ; « il existe
sur », « une » y sont tapés dans l'interligne}

(ii) Il existe une topologie sur $C'$ dont les familles couvrantes sont celles
qui sont raffinées par une famille épimorphique (ou encore, épimorphique
effective universelle ..) \emph{finie}, dans $C'$ les limites inductives finies
sont représentables, et le foncteur canonique $C \to \widetilde{C}$ y commute.

(iii) (Si dans $C$ les limites projectives finies sont représentables, de
sorte qu'on peut prendre $C' = C$.) Dans $C$ les limites inductives finies sont
représentables, et $C$ satisfait la condition ST) de I 8.8 f).\marginal{x}

(iv) Il existe un topos $E$ engendré par sa sous-catégorie pleine $E_{PF}$
(\quad), et une équivalence de catégories $C' \simeq E_{PF}$.\marginal{x}\note{les
deux « x » sont au crayon dans la marge de gauche, en regard de « I 8.8 f) »
et de la parenthèse laissée vide}

De plus, si ces conditions sont satisfaites, la topologie postulée dans (i ter)
est nécessairement celle décrite dans (ii), et le topos $E$ postulé dans (iv)
est nécessairement équivalent à $\mathrm{Ind}(C)$. Ce topos est donc parfait si
et seulement si \add{dans} $C'$ \struck{est} les limites projectives finies sont
représentables.\note{« dans » à l'encre au-dessus de la ligne, « est » tapé
et barré à l'encre}

\page{23}

\note{à l'encre bleu-noir, sur papier à en-tête de l'IHES ; les deux dernières
lignes au crayon}

SGA 4 \emph{VII} p. 20 \quad \add{algébrique, cf. \uncertain{fonction} (réf)}

\emph{Ajouter} 5.6. (a) $\widetilde{X}_{\text{ét}}$ est localement cohérent
(réf) ; il est cohérent (réf) \struck{\ill{}} \add{(réf)} \struck{\ill{}} ssi
le schéma $X$ est \struck{quasi-compact} cohérent i.e. quasi-compact
quasi-séparé (resp. quasi-séparé). \struck{\ill{}}\note{les « (réf) » sont
ajoutés au-dessus de la ligne ; « algébrique, cf. \ldots\ (réf) », sur la
première ligne, est appelé par un trait}

(b) Soit $(X_i)_{i \in I}$ système projectif filtrant à morphismes de transition
affines, $X = \varprojlim_i X_i$ (EGA IV 8), Alors
\[
\widetilde{X}_{\text{ét}} \simeq \varprojlim_{\text{top}} \widetilde{X}_i .
\]
\note{le (b) est entouré au crayon ; l'indice de $\varprojlim$ est lu « top »
(limite au sens des topos) sous réserve}

Cor. 5.6.1. $X' \in \mathrm{Ob}\,\mathrm{Et}(X)$ est (quasi-compact) dans
$X_{\text{ét}}$ \struck{\ill{}} ssi $X'$ est \struck{\ill{}} q.c. (resp. coh.)
\add{\uncertain{(quasi-sép.)}}\note{la parenthèse « (resp. coh.) » et un mot
lu « quasi-sép. » sont entourés ; le second est écrit à gauche, sous la ligne}

\ill{} \ill{} 5.7.

\note{double trait horizontal ; la suite au crayon}

SGA 4 \emph{IX} p. 42 \quad note \ill{} la p. 8

\emph{VII} 2.1 c) (p. 9) \emph{\uncertain{erroné}}

\page{24}

\note{à l'encre bleue, sans en-tête ; c'est la note annoncée à la page 23
(« IX p. 42 note \ldots\ la p. 8 »). Le cadre du haut est un ajout, sans place
marquée dans le texte}

SGA 4 \emph{IX} p. 42

\add{nous prouvons en même temps que $\widetilde{X}_{\text{ét}}$ satisfait aux
conditions équivalentes de \emph{VI} 1.25.}

(*) (Note pour la page 8.) Prouvons que la notion de constructibilité 2.3 pour
un faisceau d'ensembles $F$ coïncide avec celle de \emph{VI} 1.9.3.
\struck{Soit} La question étant \struck{en effet \ill{}} locale sur $X$, on
peut supposer $X$ affine donc « cohérent », i.e. quasi-compact et
quasi-séparé. \struck{Soit $C$ la catégorie \ill{} Soit $C$ ($C'$) la
sous-catégorie \ill{} $F$ est constructible} \struck{Il faut alors prouver que}
Soit $C$ (resp. $C'$) la sous-catégorie \add{strictement} pleine de
$\widetilde{X}_{\text{ét}}$ formée des $F$ constructibles au sens de 2.3,
\add{(resp. au sens de \emph{VI} 1.9.3), il faut prouver} \struck{que $C$ est
identique à la catégorie des objets cohérents de
$\widetilde{X}_{\text{ét}}$} \ill{} que $C = C'$.\note{les ajouts sont écrits
dans l'interligne ; plusieurs mots biffés ne se lisent pas}

\struck{En vertu de \emph{VII} 5.6 a), l'objet final $X$ de
$\widetilde{X}_{\text{ét}}$ \ill{} de $\widetilde{X}_{\text{ét}}$ (2.5). $C'$
est formé des objets \emph{cohérents} \ill{} est cohérent, donc en vertu de
\emph{VI} 1.24, et compte tenu du fait que $C$ est génératrice (2.? (ii)) dans
$\widetilde{X}_{\text{ét}}$ (2.7.2) et stable par limites inductives finies,
\ill{}. Mais il résulte de 2.7 que tout \ill{} il suffit de prouver que $C
\subset C'$, \ill{} $F \in \mathrm{Ob}\,C$ est quasi-compact, puisque un $G \in
X_{\text{ét}}$ représentable par un schéma qui est quasi-compact et
quasi-séparé est quasi-compact comme objet de $\widetilde{X}_{\text{ét}}$
(\emph{VII} 5.6.1). \ill{} $F \to F \times \ill{}$ \ill{} Comme $C$ est stable
par \struck{limites projectives} produits \ill{} (2.6 (i)) \ill{}, l'inclusion
$C \subset C'$ est donc un cas particulier de 2.1.}\note{ce passage, d'environ
seize lignes, est barré de longs traits obliques et, par endroits, de hachures
en croix ; on le donne en abrégé, avec ce qui se lit. Dans la marge de gauche,
entouré et barré : \uncertain{« \ill{} par exemple »}}

En vertu de \emph{VI} 2.9 b), il suffit de voir que $C$ est stable par limites
projectives finies, qu'elle est génératrice, \ill{} 2.6 (i) et limites
inductives finies et formée d'objets quasi-compacts. \struck{\ill{}} Compte
tenu de 2.6 (i) et (2.7.2), il reste à prouver que tout $F \in \mathrm{Ob}\,C$
est quasi-compact. Or \struck{et \emph{VI} 1.3} \struck{en vertu de} \add{en
vertu de} 2.7, \struck{puisque} \add{il existe} un $G \in
\mathrm{Ob}\,X_{\text{ét}}$ \struck{\ill{}} qui est \add{un schéma}
quasi-compact et quasi-séparé \struck{\ill{}} et un épimorphisme $G \to F$. Or
$G$ \struck{\ill{}} est un objet quasi-compact de $\widetilde{X}_{\text{ét}}$
(\emph{VII} 5.6.1), donc $F$ est quasi-compact (\emph{VI} 1.3),
cqfd.\note{l'ordre des premiers mots de cet alinéa, écrits en partie dans
l'interligne au-dessus du passage biffé, est incertain ; « 2.6 (i) » y figure
deux fois}

Bien entendu, il y aurait \add{\uncertain{notamment}} \add{eu} lieu de refondre
les §§ 1, 2 du présent exposé, \add{en les rattachant} aux notions générales de
\struck{l'exposé} finitude de \emph{VI}. Nous laissons ce soin
\struck{\ill{}} de rédaction aux générations futures.\note{« notamment » (lu
sous réserve) et « eu » au-dessus de la ligne ; sous « refondre », un mot
biffé ; « en les rattachant » appelé par un trait}

\page{25}

\note{au crayon ; le verso, dactylographié, transparaît}

SGA 4 \emph{IV} $\lbrace$ Modules plats sur anneau non commutatif, \ill{} ? ;
\uncertain{J} ! pour \struck{ensembles} \uncertain{foncteurs en groupes}
\uncertain{quelconques} \ill{} )

SGA 4 \emph{XVII} 4.2.3 \quad attention réf. à \emph{V} 3.7 (changée ?)

\emph{XVIII} \quad Théorème de finitude pour faisceaux d'anneaux de torsion
quelconque.

Préface à 2ème édition. Signaler autres contributions de
Deligne \quad | b) Intégration des torseurs ; a) Trace des torseurs sous
conditions inhabituelles

\note{trait horizontal}

Que faire avec C.D. de Verdier : inclure ?

\note{trait horizontal}

SGA 4 \emph{IV} \quad Exercices importants ; § 8 \quad | topos locaux, topos
pseudoponctuels, topos $I^{\wedge}$, $I$ un ordonné,


\page{26}
\note{feuillet dactylographié, tête-bêche, paginé « 19 » en haut au centre
(un autre nombre au crayon en haut à droite, lu \uncertain{92}) : SGA 7,
exposé VII (biextensions), 2.1.1 à 2.2.3 — la biextension comme $G_{P \times
Q}$-torseur $E$ muni de deux lois partielles, les relations (2.1.1.1) et
(2.1.1.2), les sections unité $e_{E/P}$, $e_{E/Q}$ et la section $e_E \in
\Gamma(E)$ (2.2.3). Texte publié, non reproduit. Marques : deux longues
diagonales au crayon qui annulent le feuillet (il a servi de brouillon) ; dans
la marge de gauche, au crayon, un $\psi$ entouré en regard de la ligne où la
machine a tapé le second symbole de loi ; une accolade au crayon autour de
(2.1.1.2) ; à l'encre bleue, le titre « Sections unité. » ajouté devant 2.2,
« De plus, » au-dessus d'un mot biffé devant « Ces deux trivialisations », et
« donc » inséré devant « une même section »}

\page{27}

\note{au crayon, sur une demi-feuille déchirée}

SGA 4 \emph{XVI} 1 \quad compléter comme dans Giraud

\section*{Lettre sur « 2.1.1 », et texte joint}

\page{29}

\note{lettre dactylographiée, sans date, sans formule d'appel ni signature sur
ce feuillet ; au crayon, en haut à gauche, « SGA 4 \emph{XVI} ». Les mots
tapés dans l'interligne sont donnés sans marque et signalés ; ses ajouts
manuscrits sont en \add{}. Les biffures à la machine (x) ne sont pas relevées}

VII 2

Du point de vue substance, 2.1.1 se ramène trivialement au cas où dans $X \to
S$, $X$ est affine, donc quasi-compact quasi-séparé (on dit maintenant
« cohérent »), et alors c'est une conséquence du th. de changement de base par
$f$, qui est donné dans SGA 4 XVI 1.1 et qui est essentiellement. Mais dans loc.
cit. on utilise toute la force du théorème de changement de base propre, ce qui
est en effet abusif, car en interprétant comme un énoncé d'acyclicité globale
2.1.1, on est ramené comme tu l'observes au cas particulier 2.1.5, qui se ramène
trivialement à son tour au cas du corps de degré de transcendance 1, et par là
au th. de changement de base pour le $H^1$ premier à $p$ pour une courbe propre
et lisse.\note{« et qui est essentiellement » est tapé dans l'interligne, au-dessus
de « Mais dans loc. cit. », sans suite ; « corps de » est tapé dans l'interligne
et appelé par un trait d'encre devant « base » : on le place après « au cas du »,
où il est écrit, sous réserve}
Je pense présenter les choses ainsi en réédition de SGA 4, de sorte que ton
2.1.1 y sera énoncé en toutes lettres, s'il ne l'est déjà (je n'ai pas les textes
en main). Cela te dispensera donc de le prouver, ta seule tâche consistant à
interpréter en termes de champs cet énoncé \struck{(que tu pourrais donner sous
forme de rappels)}.

Du point de vue rédaction, on a l'impression que tu t'es ingénié à tout présenter
à l'envers. Il faut en tous cas bloquer ensemble 2.1.1 et 2.1.7 (pratiquement
identiques) et 2.1.5 (celui-ci en corollaire, peut-être) \add{cf texte joint}
\struck{; ce sera dans les rappels. Avant même de donner ces énoncés, il faudra
avoir rappelé les définitions, pour que ça ait un sens.} Au lieu de pro-lisse
resp. lisse, il suffit de supposer tout \emph{localement $\underline{L}$-1-asphérique} dans
le th. \struck{(« par exemple \ldots\ »)}. D'autre part, 2.12, 2.1.3, 2.1.4,
2.1.7.1 n'ont rien à voir avec les schémas, il sont manifestement importants en
eux-mêmes et devraient passer dans un numéro à part dans \add{IV} \struck{5}
\add{ou V} (avec un peu d'ordre dans la suite des énoncés).\note{« cf texte
joint » est à l'encre au-dessus de la ligne ; le passage biffé l'est à l'encre,
encadré et barré de traits obliques. « 1- » (au-dessus de « L- », d'où la lecture « L-1- ») et « dans le th. » sont
tapés dans l'interligne et appelés à l'encre ; « il suffit de supposer tout localement
L-asphérique » est souligné à l'encre. « IV » est repassé à l'encre sur le
chiffre tapé, « 5 » barré, « ou V » écrit au-dessus}

\marginal{\ill{} \uncertain{résultats} \uncertain{évidemment} la
\uncertain{factorisation} \ill{} d'un \ill{}}\note{à l'encre, en travers de la
marge de gauche, en regard du paragraphe sur 2.12 \ldots\ 2.1.7.1}

\note{une dernière ligne, « Enfin, les compléments schématiques 2.1.7.2 à
2.1.7.4 \ldots\ », est entièrement biffée à la machine}

\page{30}

\note{le « texte joint » annoncé à la page 29 : même machine, sans pagination.
Un trait vertical au crayon longe dans la marge de gauche l'énoncé du théorème
et du corollaire. Les renvois numériques entre parenthèses (1.11, 2.1, 1.13, (ii))
sont écrits à la main dans des blancs de la machine}

\emph{Théorème} 2.1.1. Soit $f : S' \to S$ un morphisme universellement
localement $\underline{L}$-1-asphérique (SGA 4 XV \add{1.11}), par exemple
(SGA \add{4} XV \add{2.1}) un morphisme lisse, ou tel que \add{$S'$} soit
isomorphe à la limite projective d'un système projectif filtrant de
$S$-schémas lisses $S'_i$ à morphismes de transition affines (SGA 4 XV
\add{1.13 (ii)}). Soit en $g : X \to S$ un morphisme de schémas, d'où un carré
cartésien\note{« projective » et « filtrant » sont tapés dans l'interligne et
appelés à l'encre ; « $S'$ » est écrit à l'encre sur une lettre tapée ; « en »
est tapé au-dessus de « Soit », sans correction de la suite}

\begin{tikzcd}
X \arrow[d, "g"'] & X' \arrow[l, "f'"'] \arrow[d, "g'"] \\
S & S' \arrow[l, "f"]
\end{tikzcd}

\add{(\emph{V} 5 \quad)}\note{à l'encre sous le carré, avec deux petits traits
contre la flèche $g$}

$G$ un \add{ind-$\underline{L}$-}champ sur $X$, $G' = f'^{*}(G)$ le champ image
inverse sur $X'$.\note{« ind-L- » est écrit à l'encre au-dessus de « champ » ; la
suite tapée de la ligne, « dont les faisceaux d'automorphismes sont des
ind-L-faisceaux », est biffée à la machine}

a) Si $f$ est un morphisme \emph{local} de schémas strictement locaux, alors le
foncteur
\[
G(X) \longrightarrow G'(X')
\]
est une équivalence de catégories.

b) Si $g$ est cohérent (i.e. quasi-compact et quasi-séparé) alors le morphisme
de champs « de changement de base »
\[
f^{*}(g_{*}(G)) \longrightarrow g'_{*}(f'^{*}(G))
\]
est une équivalence de champs sur $S'$.\note{les étoiles des foncteurs sont à la
main ; un mot tapé après la formule est biffé à l'encre}

\emph{Corollaire} 2.1.2. Soient $k$ un corps séparablement clos, $k'$ une
extension séparablement close de $k$, $X$ un $k$-schéma, $X' = X \otimes_k k'$,
alors pour tout ind-$\underline{L}$-champ $G$ sur $X$, désignant par $G'$ son
image inverse sur $X'$, le foncteur
\[
G(X) \longrightarrow G'(X')
\]
est une équivalence de catégories.

On sait en effet (SGA 4 XVI \quad) que $\mathrm{Spec}(k') \to \mathrm{Spec}(k)$
est universellement localement $\underline{L}$-1-asphérique.

Démonstration. On prouve a) par référence à SGA 4 XVI (nouvelle édition !)
\add{et à \emph{V} ---}, et b) en se ramenant à a) par des lemmes 2.1.7.2 et
2.1.7.4 renumérotés. \add{En tout, ça devrait faire 3 pages au Chap. \emph{VII}
au maximum, --- mais il y aura \ill{} \ill{} au Chap. \emph{I}.}\note{la
dernière phrase est écrite à l'encre au bout de la ligne et sous elle}

\section*{Notes pour la réédition de SGA 4 (suite)}

\page{32}

\note{sur papier à en-tête de l'IHES ; encre bleu-noir, une ligne à l'encre
rouge, la dernière au crayon. Les points repris de la lettre d'Artin (pages
33-34) : XIII 3.1, XIX p. 7 et 18, XII 5.10, IX p. 34}

\emph{Introduction} \uncertain{vider} \quad | Bibliographie \ill{}, mention de
C.T.

\emph{Artin} SGA 4

\uncertain{réf} $\times$ [1]

$\times$ p. 7 \quad Début de phrase manque

\note{trait horizontal}

\emph{XII} \quad compléter réf : Giraud

5.10. \quad Compléter \ill{} par réf. aux travaux récents d'Artin

\emph{IX} p. 34 \quad réf. th. de Tsen\note{cette ligne à l'encre rouge}

\note{trait horizontal}

\emph{XIII} \struck{\ill{}} 3.1. \quad \uncertain{Supprimer} $\times$
\uncertain{noethérien} (\ill{}, p. 13 l. $-3$ \ill{} \ill{})

\note{trait horizontal ; la suite au crayon}

XIX \quad p. 3, 9, 10 \qquad p. 7, \struck{\ill{}}, 18, 55, 39, 43

\page{33}
\note{lettre dactylographiée en anglais, sans date ni signature sur ce feuillet,
de M. Artin (voir la page 34) : réponses à ses questions sur SGA 4. Sur XIII 3.1,
Artin accepte l'hypothèse noethérienne si Grothendieck la juge nécessaire ; sur
XIX, il accepte de remplacer $\mu_n$ par $\mu_n(A)$ et la suggestion pour p. 7
ligne 2, mais pense que p. 18 ligne $-4$ doit rester (le cas $q = 0$ demande un
traitement particulier) et devra réfléchir au reste. Non reproduit. Marques :
un trait rouge en travers du paragraphe sur XIII 3.1 ; le paragraphe sur XIX
entouré d'un cadre à l'encre rouge ; deux traits horizontaux au crayon qui
séparent les paragraphes}

\page{34}
\note{suite de la même lettre : la référence [1] pour l'exposé X (J. Ax, « Some
conjectures of Serre on cohomological dimension », Proc. AMS 16, 1965) ; pour X
p. 7, un paragraphe en français à ajouter (le corps résiduel de $R'$ est $k$,
les racines $\ell$-ièmes de 1 sont dans $K$, et un diagramme $H' \hookrightarrow
G' \to \overline{G}$, $H \hookrightarrow G \to \overline{G}$) ; pour XII 5.10, la
référence à « my paper on alg. approx », M. Artin, « Algebraic approximation of
structures over complete local rings » (à paraître), théorème (3.1) ; pour IX
p. 34, des références pour le théorème de Tsen (Douady, Séminaire Bourbaki
exposé 189 ; Lang, « Rapport sur la cohomologie des groupes », 1966 ; Tsen,
1933). Non reproduit. Les $\ell$ tracés à la main sont ceux de la machine.
Marques : quatre longs traits obliques au crayon et un à l'encre rouge qui
traversent le feuillet, deux traits horizontaux au crayon entre les
paragraphes. Rien d'écrit de sa main}

\page{35}

\note{à l'encre bleue ; les deux premiers alinéas sont annulés au crayon, le
premier par trois traits obliques, le deuxième et le début du troisième par
de grandes boucles}

\struck{Exposés \emph{XI} à \emph{XIV} sont tapés avec des caractères plus gros
que le reste du Séminaire.}

\struck{Vider \ill{} (des exposés \emph{XVI}, \emph{XIV} \ldots) ?}

\struck{\emph{VI} \quad p. 49, p. 50 \quad \uncertain{univers} est \emph{non
vide} ! ; p. 52 \quad vider \uncertain{footnote}}

\emph{VIII} \quad généralités sur foncteurs fibres : renvoyer à \struck{SG}
\emph{IV} \ldots

\page{37}

\note{au crayon, sur papier à en-tête de l'IHES}

SGA 4 \emph{XVIII} \quad (\uncertain{Deligne})\note{le nom est entouré, et
relié par un trait à « XVIII »}

Références pour SGA 5 \emph{VII} \quad $\lbrace$ Dém. de 3.5.3 ; 1.1 b)

\note{trait horizontal}

\emph{V} \quad dire quelques mots sur relations avec $\varprojlim$ et
$\varprojlim^{(p)}$ :

a) Sur un topos de la forme $\widehat{C}$, les foncteurs $H^n$ sont les
$\varprojlim^{(n)}_{C}$ \ill{}

b) Je \ill{} plus \ill{} telles interprétations pour un topos général,
\ill{} \uncertain{puisqu'on} est obligé de prendre \ill{}
\[
H^n(\widetilde{E}, F) \simeq \varinjlim_{\mathcal{U} \in R}\
\varprojlim_{\mathcal{U}} \ill{}
\]
\marginal{Si ça marche \ldots}\note{le $E$ est écrit sur un $C$ ; après
$\varprojlim_{\mathcal{U}}$, un groupe de signes repassé (« $F$ \ldots\ / \ldots »)
n'est pas lu. La note de droite est séparée par un trait vertical}

$\mathcal{U}$ crible variable

On a un recouvrement de l'objet final de $\widehat{C}$ par les objets de $C$

$X \times Y$\note{très pâle, au bas de la feuille}

\section*{SGA 4 VIII, § 10}

\page{38}

\note{tapuscrit coté au crayon en haut à gauche « SGA 4 \emph{VIII} » (le chiffre
romain surligné) ; section non publiée (voir l'en-tête), donnée entière. Les
numéros de renvoi à IV (11, 9) sont écrits au crayon dans des blancs de la
machine ; les « x » au crayon de la marge de gauche marquent ces lignes. Les
biffures à la machine ne sont pas relevées}

10. \quad \emph{Compléments} : \emph{espaces algébriques et étendues
algébriques}.

Nous donnons sous forme d'exercices quelques compléments, qui ne seront pas
utilisés dans la suite du Séminaire.

10.1. \quad Soit $(X, \underline{O}_X)$ un schéma, \struck{posons
$X_{\text{ét}} = \mathrm{Top}(X)$, de} sorte que $(\widetilde{X}_{\text{ét}},
\underline{O}_X)$ est un topos annelé. Montrer qu'on peut reconstruire le
schéma $X$ à isomorphique unique près, quand on connait le topos annelé
$(\widetilde{X}_{\text{et}}, \underline{O}_X)$ à équivalence près. (Utiliser 6.1
pour reconstituer le topos $\mathrm{Top}(X) =
\mathrm{Ouv}(\widetilde{X}_{\text{ét}})^{\sim}$, puis utiliser IV
4.\add{2.3}.) Plus précisément, montrer que si $X$, $X'$ sont deux schémas,
alors l'application naturelle
\[
\mathrm{Hom}(X, X') \longrightarrow
\mathrm{Homtoplocan}(\widetilde{X}_{\text{ét}}, \widetilde{X}'_{\text{ét}})
\]
(notations de IV \add{11}. \quad) induit une équivalence entre la catégorie
discrète définie par l'ensemble $\mathrm{Hom}(X, X')$ et la catégorie
\emph{Homtoplocan}$(\widetilde{X}_{\text{ét}}, \widetilde{X}'_{\text{ét}})$ (où pour
simplifier on désigne par $\widetilde{X}_{\text{ét}}$ le topos étale de $X$
considéré comme topos annelé). (Utiliser IV 4.\add{2.3}.)\note{« posons
\ldots\ de » est barré à l'encre bleue ; dans les deux renvois IV 4.2.3, le
« 2 » est à l'encre bleue et le « 3 » au crayon. Un signe au crayon dans la
marge de droite, en regard de la seconde parenthèse, n'est pas lu}

10.2. \quad Soit $(E, \underline{O}_E)$ un topos annelé. Montrer que les
conditions suivantes sont équivalentes (comparer IV \add{9}. \quad) :

(i) Il existe une famille $(X_i)_{i \in I}$ d'objets de $E$ couvrant l'objet
final, et une famille $(X'_i)_{i \in I}$ de schémas, tels que pour tout $i \in
I$, le topos annelé induit $E_{/X_i}$ soit équivalent à
$\widetilde{X'_i}_{\text{ét}}$.

(ii) Comme (i), avec $I$ réduit à un élément.

(iii) Il existe un schéma $X_0$, et une prérelation d'équivalence \add{$X_1$}
dans $X_0$ au sens de la catégorie (Sch) des schémas, \struck{\ill{}}, telle
que $R$ soit étale i.e. le morphisme $p_1 : \add{X_1} \to X_0$ est étale, et que
le topos annelé $(E, \underline{O}_E)$ soit équivalent au topos annelé déduit
par recollement du diagramme \add{\uncertain{$\frac{1}{2}$-simplicial}} de topos
annelés \add{étales associés au diagramme de schémas}
\[
X_2 \mathrel{\substack{\longrightarrow\\[-2pt] \longrightarrow\\[-2pt]
\longrightarrow}} X_1 \rightrightarrows X_0
\]
\note{les deux $X_1$ sont écrits à l'encre sur un $R$ tapé, et l'indice 0 de
$X_0$ est repassé ; une parenthèse tapée est noircie à l'encre bleue. Les
ajouts « \ldots-simplicial » (au-dessus de « de topos annelés »), « étales » et
« associés au diagramme de schémas » sont à l'encre bleue ; le premier, petit,
est lu « ½-simplicial » (= semi-simplicial) sous réserve. Les trois flèches
et les indices du diagramme sont tracés à la main}

\page{39}

\note{suite du même tapuscrit}

où $X_2 = (X_1, p_2) \times_{X_0} (X_1, p_1)$ (cf. IV \add{9} \quad pour le
procédé de recollement de topos par données de descente).

Lorsque les conditions précédentes sont satisfaites, on dira que le topos
annelé $(E, \underline{O}_E)$ est une \emph{étendue algébrique} (comparer avec
la notion d'étendue topologique IV \add{9}. \quad).

10.3. \quad Montrer qu'une étendue algébrique a suffisamment de points. Une
étendue algébrique est \emph{localement} annelée (IV \add{11}. \quad), et si
$E$, $E'$ sont deux étendues algébriques, la catégorie
\emph{Homtoplocan}$(E, E')$ (IV \add{11}. \quad) est un groupoïde (toute
flêche est inversible). (Procéder pour le premier point comme dans IV \add{9}.
\quad ; pour le deuxième, utiliser IV \quad .)\note{deux « x » au crayon dans
la marge en regard de cette dernière ligne}

10.4. \quad \add{(Comparer \emph{IV} 9. \quad)} Soit $E$ une étendue algébrique.
Montrer que les conditions suivantes sont équivalentes :\note{la parenthèse est
à l'encre bleue au-dessus de la ligne, avec un « x » au crayon dans la marge}

(i) Pour toute étendue algébrique $E'$, la catégorie
\emph{Homtoplocan}$(E', E)$ est rigide, i.e. (comme c'est un groupoïde) elle
est équivalente à une catégorie discrète.

(ii) Même énoncé, $E'$ étant le topos ponctuel annelé par un corps
algébriquement clos. En d'autres termes, pour tout point géométrique (IV
\add{11}) de $E$, le groupe des automorphismes de $p$ est le groupe unité.

(iii) (Lorsque $E$ est exprimé en termes d'une prérelation d'équivalence étale
$X_1$ dans un schéma $X_0$, comme dans 10.2 (iii).) La prérelation
d'équivalence $(X_1, p_1, p_2)$ dans le schéma $X_0$ est une relation
d'équivalence, i.e. le morphisme $(p_1, p_2) : X_1 \to X_0 \times X_0$ est un
monomorphisme.\note{« étale » est tapé dans l'interligne et appelé à l'encre
bleue}

Lorsque ces conditions sont satisfaites, on dit que l'étendue algébrique $E$
est \emph{rigide}, ou encore que $E$ est un \emph{espace algébrique}
\add{(*)}.

10.5.\note{numéro tapé seul, entouré au crayon ; le reste du feuillet est vide
jusqu'à la note}

\add{(*) La terminologie est due à M. ARTIN, qui a mis en évidence l'importance
de cette notion en Géométrie Algébrique.}\note{note de bas de page écrite à la
main, à l'encre, appelée par l'astérisque ajouté après « espace algébrique »}

\section*{Terminologie, notations SGA 4}

\page{40}

\add{Terminologie, notations SGA 4 Vol. 2 \quad Attention, faire \emph{deux}
listes disjointes !}\note{à l'encre bleue, en deux lignes, dans un cadre ouvert
en haut du feuillet ; le « 2 » est repassé. La liste, dactylographiée, donne
pour chaque entrée le terme français, le terme anglais et le renvoi (exposé et
numéro de SGA 4), séparés ici par des tirets, comme dans la suite (pages 42-44,
lot 3). Les notations, encore mêlées aux termes, sont entourées au crayon pour
la seconde liste : on le signale par « (entouré) ». Les biffures à la machine
ne sont pas relevées}

\begin{itemize}
\item (entouré) (Sch) --- VII 1.1 ; $\mathrm{Et}_{/X}$, $X_{\text{ét}}$ ---
VII 1.2\note{un seul ovale au crayon entoure les deux lignes}
\item Topologie étale (sur (Sch), sur un schéma) --- étale topology (on (Sch),
on a scheme) --- VII 1.2\note{un trait d'encre bleue relie « on (Sch), » à la
fin de l'anglais, tapée à la ligne suivante}
\item site étale d'un schéma --- étale site of a scheme --- VII 1.2
\item topos étale d'un schéma --- étale topos of a scheme --- VII 1.2
\item (entouré) $f_{\text{ét}}$ --- VII \uncertain{1.4}\note{le numéro est
surchargé à l'encre bleue (un « 4 » sur le dernier chiffre) ; lecture sous
réserve}
\item point géométrique d'un schéma --- geometric point of a scheme --- VIII 3.1.
\item foncteur fibre (relatif à un point géométrique d'un schéma) --- fiber
functor (corresponding to a geometric point of a scheme) --- VIII 3.3\note{une
lettre tapée après « functor » est barrée à l'encre bleue}
\item (entouré) $F_{\xi}$ --- VIII 3.3\note{le $\xi$ est tracé à la main}
\item hensélien (anneau local ---) --- henselian local ring --- VIII 4.1
\item hensélisé d'un anneau local --- henselization of a local ring --- VIII 4.1
\item strictement local (anneau) --- strictly local ring --- VIII 4.2
\item --- id. --- (schéma) --- strictly local scheme --- VIII 4.2
\item anneau strictement local (d'un schéma en un point géométrique) ---
strictly local ring (of a scheme at a geometric point) --- VIII 4.3
\item hensélisé strict d'un anneau local --- strict henselisation of a local
ring --- VIII 4.4
\item localisé strict d'un schéma en un point géométrique --- strict
localization of a scheme at a geometric point --- VIII 4.3
\item flêche de spécialisation --- specialization arrow --- VIII 7.2\note{« arrow »
est tapé au-dessus d'un mot biffé à la machine}
\item spécialisation d'un point géométrique --- specialization of a geometric
point --- VIII 7.2\note{un petit trait d'encre sépare « VIII » de « 7.2 », ici
et à la ligne suivante}
\item générisation d'un point géométrique --- generization of a geometric point
--- VIII 7.2
\item (entouré) $\overline{X}(\xi)$ ($\xi$ point géométrique du schéma $X$) ---
VIII 7.1\note{les deux $\xi$ sont tracés à l'encre bleue}
\item Homomorphisme de spécialisation --- specialization homomorphism --- VIII
7.7
\item (entouré) $\underline{\mathrm{Pt}}(E)$ --- VIII 7.8.
\end{itemize}

\end{document}
